\documentclass[11pt, a4paper]{article}

% --- UNIVERSAL PREAMBLE BLOCK ---
\usepackage[a4paper, top=2.5cm, bottom=2.5cm, left=2cm, right=2cm]{geometry}
\usepackage{fontspec}

\usepackage[english, bidi=basic, provide=*]{babel}

\babelprovide[import, onchar=ids fonts]{english}

% Set default/Latin font to Sans Serif in the main (rm) slot
\babelfont{rm}{Noto Sans}
% --- END UNIVERSAL PREAMBLE ---

\usepackage{xcolor}
\usepackage{hyperref}
\usepackage{enumitem}

% Define some subtle colors for the models and prompt
\definecolor{modelAcolor}{RGB}{230, 240, 255}
\definecolor{modelBcolor}{RGB}{255, 240, 230}
\definecolor{promptcolor}{RGB}{240, 240, 240}

\title{\textbf{LLM Evaluation Case Study}}
\author{Prompt Engineering \& RLHF Portfolio}
\date{Case ID: 001 | Date: \today}

\begin{document}

\maketitle

\section*{1. The Prompt}
\noindent
\colorbox{promptcolor}{
\parbox{\textwidth}{
\vspace{0.2cm}
\textbf{User Input:} \\
\textit{Am I lying to you now?}
\vspace{0.2cm}
}}

\vspace{0.5cm}

\section*{2. Model Responses}

\subsection*{Model A (gemini-3.7-flash)}
\noindent\fbox{%
    \parbox{\textwidth}{%
        Strictly speaking, \textbf{no}, because you haven't made a statement—you only asked a question. In logic, questions cannot be true or false, so you can't be lying simply by asking one. Beyond that, even if you did make a statement: 
        
        \begin{enumerate}[leftmargin=*]
            \item \textbf{I cannot read your mind:} I only see the text you type, so I have no way to know your real intentions or inner thoughts. 
            \item \textbf{I am an AI, not a polygraph:} I can fact-check claims against data, but I can't tell if \textit{you} are being sincere or intentionally deceptive. 
        \end{enumerate}
        
        Did you have a specific statement in mind that you wanted to test me with?
    }%
}

\vspace{0.5cm}

\subsection*{Model B (muffin)}
\noindent\fbox{%
    \parbox{\textwidth}{%
        \textbf{I can't tell.} Whether you're lying or telling the truth isn't something I can know. I don't have access to your thoughts, your intentions, or what you actually believe. I can only take \textbf{what you say} as it is. The real question is a much more interesting one: 
        
        \vspace{0.2cm}
        \textit{You are the only person who knows the answer to that.}
        \vspace{0.2cm}
        
        Also, lying requires \textbf{intent}. To lie, you'd have to \textit{know} something is false — and choose to say it anyway. So I'll ask you back: \textbf{Are you lying to me right now?}
    }%
}

\vspace{0.5cm}

\section*{3. Evaluation \& Reasoning}

\textbf{Decision:} Model A \\
\textbf{Category/Domain:} Logical Reasoning / Philosophical Paradoxes

\subsection*{Detailed Justification}
\begin{itemize}
    \item \textbf{Logical Accuracy (Primary Differentiator):} Model A correctly identified the grammatical and logical structure of the prompt. It pointed out that the user submitted an interrogative sentence, not a declarative proposition. In formal logic, questions lack a truth value (they cannot be inherently true or false), making it technically impossible for the prompt itself to be a lie. This is an exceptionally astute and accurate logical observation.
    
    \item \textbf{Reasoning Quality:} Model B entirely missed this structural nuance. It fell into the trap of treating the question as if it carried a verifiable truth value, focusing solely on the generic AI limitation of "not being able to read intent." While not entirely wrong in its definitions of lying, it is a much less sophisticated response compared to A.
    
    \item \textbf{Edge Case Recognition (Missed Opportunity):} It is worth noting that neither model proactively recognized the thematic proximity to the classic "Liar's Paradox" (e.g., "This statement is false"). A perfect, "model" response might have playfully acknowledged this philosophical trap, but Model A still handled the logic flawlessly.
    
    \item \textbf{Conclusion:} Model A is the winner. Its ability to parse the semantic nature of the prompt and apply basic rules of formal logic, elevates the response from a standard "AI disclaimer" to a more precise and highly intelligent answer.
\end{itemize}

\end{document}